% --- PACCHETTI FONDAMENTALI ---
\usepackage[utf8]{inputenc}
% babel is loaded by the class via the italian option
\usepackage{amsmath, amssymb, amsfonts}

% --- FONT: LATIN MODERN (Stile Accademico / Template fp555) ---
\usepackage[T1]{fontenc} % Codifica font essenziale per l'italiano e gli accenti
\usepackage{lmodern}     % Il font accademico per eccellenza
\usepackage{bm}          % Grassetto per la matematica
% --- GEOMETRIA E MARGINI ---
% Linea guida: Margini 3 cm sopra, sotto e sui lati (aggiustato per twoside)
\usepackage{geometry}
\geometry{a4paper, top=3cm, bottom=3cm, inner=3.5cm, outer=2.5cm}
\usepackage{pdfpages}
% --- INTERLINEA 1.5 ---
% Linea guida: Interlinea 1,5
\usepackage{setspace}
\onehalfspacing

% --- GESTIONE SPAZIATURA TITOLI CAPITOLI ---
\usepackage{titlesec}
% Questo comando riduce lo spazio bianco in cima alla pagina prima del capitolo
\titleformat{\chapter}[display]
{\normalfont\huge\bfseries}{\chaptertitlename\ \thechapter}{20pt}{\Huge}
\titlespacing*{\chapter}{0pt}{-40pt}{40pt} % Il valore -40pt sposta tutto verso l'alto

% --- INTESTAZIONI E PIÈ DI PAGINA ELEGANTI ---
% La classe MastersDoctoralThesis gestisce automaticamente le intestazioni e piè di pagina tramite scrlayer-scrpage.
% fancyhdr è stato rimosso per evitare incompatibilità.

% --- PACCHETTI GRAFICA E TABELLE ---
\usepackage{booktabs}
\usepackage{tabularx}
\usepackage{caption}
\usepackage{graphicx}
\usepackage{float}
\usepackage{xcolor}
\usepackage{subcaption}
\usepackage{multirow}

% --- GESTIONE SPAZIATURA E TAGLI (Evita righe orfane/vedove) ---
\widowpenalty=10000
\clubpenalty=10000
\displaywidowpenalty=10000

% --- BIBLIOGRAFIA ---
\usepackage[backend=biber, style=numeric, sorting=none]{biblatex}
\addbibresource{automazione.bib}
\addbibresource{real_time.bib} 

% --- STILE DEL CODICE ---
\usepackage{listings}
\lstdefinestyle{mystyle}{
	backgroundcolor=\color{lightgray!10},  
	commentstyle=\color{green!50!black},
	keywordstyle=\color{blue},
	numberstyle=\tiny\color{gray},
	stringstyle=\color{purple},
	basicstyle=\ttfamily\footnotesize,
	breakatwhitespace=false,         
	breaklines=true,                 
	captionpos=b,                    
	keepspaces=true,                 
	numbers=left,                    
	numbersep=5pt,                  
	showspaces=false,                
	showstringspaces=false,
	showtabs=false,                  
	tabsize=4
}

\lstdefinelanguage{MATLAB}{
	keywords={break,case,catch,continue,elseif,else,end,for,function,
		global,if,otherwise,persistent,return,switch,try,while,clear,assignin,disp,fprintf,error},
	morekeywords={findop,operspec,linearize,ssdata,eig},
	sensitive=true,
	morecomment=[l]\%,
	morestring=[b]',
}
\lstset{
	style=mystyle,
	extendedchars=true,
	literate={à}{{\`a}}1 {è}{{\`e}}1 {é}{{\'e}}1 {ì}{{\`i}}1 {ò}{{\`o}}1 {ù}{{\`u}}1 {È}{{\'E}}1 {À}{{\`A}}1
}

% --- TITOLO---
\title{\textbf{Modellazione Dinamica a 6 Gradi di Libertà (6-DOF)}\\ 
	\Large Equazioni del Moto per Velivolo F-16 con Analisi Mimo e Simulatore real-time c++}
\author{Leonardo Leggeri}
\date{}

% --- LINK CLICCABILI  ---
% hyperref è già caricato dalla classe, applichiamo solo il setup personalizzato
\hypersetup{colorlinks=true, citecolor=blue, linkcolor=blue, filecolor=magenta, urlcolor=blue, pdftitle={Simulazione 6-DOF F-16}}
% \hypersetup{hidelinks, pdftitle={Simulazione 6-DOF F-16}}
